% Discrete three-type version of Ide & Talamas (2025), stated and proved in full.
% No numerical examples: every claim is a theorem about the primitives.
% Compile with: latexmk -pdf discrete-three-type.tex
\documentclass[11pt]{article}
\usepackage[utf8]{inputenc}
\usepackage[T1]{fontenc}
\usepackage{lmodern}
\usepackage[margin=2.4cm]{geometry}
\usepackage{amsmath,amssymb,amsthm}
\usepackage{enumitem}

\theoremstyle{plain}
\newtheorem{teorema}{Teorema}
\newtheorem{lema}{Lema}
\newtheorem{corolario}{Corolario}
\theoremstyle{definition}
\newtheorem{definicion}{Definición}
\theoremstyle{remark}
\newtheorem*{intuicion}{Intuición}
\newtheorem*{obs}{Observación}
\renewcommand{\proofname}{Demostración}

\newcommand{\ai}{a}
\newcommand{\N}{\mathsf{N}}
\newcommand{\Z}{\mathsf{Z}}
\newcommand{\Mk}{\mathsf{M}}

\title{Versión discreta de tres tipos\\[4pt]
\large Ide \& Talamàs (2025): las Proposiciones 5 y 6 demostradas con tipos finitos}
\author{Alejandro Ventura --- \texttt{ai-05-ide}}
\date{Fuente: arXiv:2312.05481v11}

\begin{document}
\maketitle

\begin{abstract}
\noindent
Construimos una versión de tipos finitos del modelo de Ide y Talamàs (2025) y demostramos, sin
ejemplos numéricos y sin condiciones auxiliares sobre los primitivos, que: (i) los salarios de
equilibrio son únicos y vienen dados por un máximo sobre configuraciones de firma; (ii) hay ganadores
en el fondo de la distribución si y solo si la capacidad de la IA supera un umbral explícito;
(iii) ese umbral es estrictamente mayor con IA autónoma que sin autonomía, en un factor igual al
inverso del costo de oportunidad de la IA; y (iv) los más conocedores prefieren siempre la IA
autónoma, pero pierden respecto del mundo sin IA. El punto (iv) es donde la versión discreta se
separa del paper, y demostramos exactamente por qué.
\end{abstract}

\section{Entorno}

\paragraph{Primitivas.} Tres tipos de humanos con conocimiento $z_1<z_2<z_3$ en $[0,1)$ y masas
$\mu_1,\mu_2,\mu_3>0$. Cada humano tiene una unidad de tiempo. Cada oportunidad de producción trae un
problema de dificultad $x\sim U[0,1]$: quien lo enfrenta produce una unidad si su conocimiento excede
$x$. Las oportunidades son abundantes. Cada consulta consume $h\in(0,1)$ unidades del tiempo del
solucionador, de modo que una firma de dos capas con trabajadores de conocimiento $z$ contrata
\[
  n(z)=\frac{1}{h(1-z)}>1
\]
de ellos por solucionador. Escribimos $n_i\equiv n(z_i)$ y $n_a\equiv n(\ai)$. La IA convierte compute
en agentes de capacidad $\ai\in[0,z_3)$, sustitutos perfectos de un humano con ese conocimiento; el
stock de compute es $\mu$ y su alquiler es $r$. El precio del producto está normalizado a 1.

\paragraph{Configuraciones de firma.} Con tres tipos y una IA el conjunto de configuraciones es
finito. Sus beneficios son
\begin{align*}
  \Pi^{1}_i &= z_i-w_i, &&\text{una capa, humano } i;\\
  \Pi^{2}_{i,j} &= n_j\,[z_i-w_j]-w_i, &&\text{solucionador } i,\ \text{trabajadores } j,\ z_j\le z_i;\\
  \Pi^{1}_{AI} &= \ai-r, &&\text{una capa, IA};\\
  \Pi^{tA}_{j} &= n_j\,[\ai-w_j]-r, &&\text{solucionador IA, trabajadores } j,\ z_j\le \ai;\\
  \Pi^{bA}_{i} &= n_a\,[z_i-r]-w_i, &&\text{solucionador } i,\ \text{trabajadores IA},\ z_i\ge \ai.
\end{align*}
Cuando la IA es \emph{no autónoma} no puede perseguir oportunidades de producción, de modo que
$\Pi^{1}_{AI}$ y $\Pi^{bA}_{i}$ desaparecen del conjunto admisible.

\begin{definicion}[Equilibrio competitivo]\label{def:eq}
Un equilibrio es un vector $(w_1,w_2,w_3)\in\mathbb R_{\ge0}^{3}$, un alquiler $r\ge0$ y una
asignación de los humanos y del compute a configuraciones admisibles tales que
\begin{description}[leftmargin=3em]
  \item[$(\N)$] ninguna configuración admisible tiene beneficio estrictamente positivo;
  \item[$(\Z)$] toda configuración usada en la asignación tiene beneficio exactamente cero;
  \item[$(\Mk)$] todos los humanos están empleados, el compute usado no excede $\mu$ y, si queda
  compute ocioso, $r=0$.
\end{description}
\end{definicion}

\begin{lema}[las configuraciones puras bastan]\label{lem:puras}
Considérese una firma de dos capas cuyo solucionador ---humano del tipo $i$ o agente de IA--- dedica
una fracción $t_j\ge0$ de su tiempo a supervisar trabajadores del tipo $j$, con $\sum_j t_j\le1$, y
emplea el tiempo restante en producir por su cuenta. Su beneficio no excede al máximo de los
beneficios de las configuraciones puras listadas arriba. En consecuencia, si ninguna configuración
pura tiene beneficio positivo, tampoco lo tiene ninguna firma mixta.
\end{lema}
\begin{proof}
Un trabajador del tipo $j$ consume $h(1-z_j)=1/n_j$ unidades del tiempo del solucionador, así que con
una fracción $t_j$ de tiempo este supervisa $t_j n_j$ trabajadores de ese tipo. El beneficio de una
firma con solucionador humano del tipo $i$ es entonces
\[
  \Pi(t)=\sum_j t_j\,n_j\,[z_i-w_j]+\Big(1-\sum_j t_j\Big)z_i-w_i ,
\]
y el de una firma con solucionador de IA,
$\Pi(t)=\sum_j t_j\,n_j\,[\ai-w_j]+\big(1-\sum_j t_j\big)\ai-r$. En ambos casos $\Pi$ es afín en
$t$ sobre el politopo compacto y convexo $\Delta=\{t\ge0:\ \sum_j t_j\le1\}$, de modo que alcanza su
máximo en un punto extremo. Los puntos extremos de $\Delta$ son el origen y los vectores unitarios
$e_j$: el origen da $\Pi^{1}_i$ (o $\Pi^{1}_{AI}$) y $e_j$ da $\Pi^{2}_{i,j}$ (o $\Pi^{tA}_{j}$).
Las firmas que mezclan trabajadores humanos con agentes de IA se tratan igual, añadiendo una
coordenada cuyo vértice es $\Pi^{bA}_{i}$.
\end{proof}

\begin{intuicion}
Mezclar no sirve porque el único recurso escaso de la firma es el tiempo del solucionador, y cada
tipo de trabajador le rinde un excedente \emph{por unidad de tiempo} constante: conviene dedicar todo
el tiempo al mejor de ellos. Esto es lo que permite enumerar un conjunto finito de configuraciones sin
pérdida de generalidad; sin este lema, la condición $(\N)$ verificaría solo una parte de las
desviaciones posibles.
\end{intuicion}

\paragraph{Supuestos.} Escribimos $\bar z\equiv\dfrac{z_1}{1-h(1-z_1)}$, bien definido porque
$h(1-z_1)<1$.
\begin{description}[leftmargin=3em]
  \item[(A0)] $z_1\le \ai<z_3$;
  \item[(A1)] $\mu_1>n_1(\mu_2+\mu_3)$;
  \item[(A2)] $\mu_2>n_2\,\mu_3$;
  \item[(A3)] $\mu>\dfrac{\mu_1}{n_1}+\dfrac{\mu_2}{n_2}+n_a(\mu_2+\mu_3)$.
\end{description}

\begin{intuicion}
(A0) dice que la IA es al menos tan capaz como el tipo menos conocedor y no alcanza al más
conocedor; sin la primera mitad, el tipo 1 podría supervisar agentes de IA y la configuración
$\Pi^{bA}_{1}$ entraría en escena. (A1) y (A2) dicen que cada tipo es demasiado numeroso para caber entero bajo los solucionadores
humanos que están por encima de él: siempre sobra gente del tipo 1 y del tipo 2. (A3) dice que el
compute alcanza para todos los usos posibles y todavía sobra. Son los análogos discretos de dos
supuestos del paper: que existen productores independientes y que el compute es abundante respecto
del tiempo humano.
\end{intuicion}

\section{Cuatro lemas}

\begin{lema}[cotas inferiores]\label{lem:cotas}
En cualquier equilibrio, para todo $i$ y todo $j$ con $z_j\le z_i$:
\[
  w_i\ge z_i,\qquad w_i\ge n_j\,(z_i-w_j),\qquad
  w_j\ge \ai-\frac{r}{n_j}\ \ (z_j\le \ai),\qquad
  w_i\ge n_a\,(z_i-r)\ \ (z_i\ge \ai),
\]
donde las dos últimas solo aplican cuando la configuración correspondiente es admisible.
\end{lema}
\begin{proof}
Son las cuatro desigualdades $\Pi^{1}_i\le0$, $\Pi^{2}_{i,j}\le0$, $\Pi^{tA}_j\le0$ y
$\Pi^{bA}_i\le0$ de la condición $(\N)$, despejando el salario.
\end{proof}

\begin{lema}[el salario es un valor alcanzado]\label{lem:alcanzado}
En cualquier equilibrio, el salario de cada tipo coincide con el valor de cero beneficio de alguna
configuración admisible que lo emplea.
\end{lema}
\begin{proof}
Por $(\Mk)$ cada tipo está empleado en alguna configuración de la asignación; por $(\Z)$ esa
configuración tiene beneficio cero; despejando el salario del tipo en cuestión se obtiene el valor
afirmado.
\end{proof}

\begin{lema}[las masas excluyen una rama]\label{lem:masas}
Bajo (A1) ningún equilibrio tiene a todos los individuos del tipo 1 empleados como trabajadores de
solucionadores humanos. Bajo (A2) ocurre lo mismo con el tipo 2 respecto de los solucionadores del
tipo 3. Además, en ningún equilibrio el tipo 3 es trabajador.
\end{lema}
\begin{proof}
La capacidad total de los solucionadores humanos por encima del tipo 1 es $n_1(\mu_2+\mu_3)<\mu_1$ por
(A1), y la de los del tipo 3 sobre trabajadores del tipo 2 es $n_2\mu_3<\mu_2$ por (A2); en ambos
casos no alcanza para emplear a toda la masa. Para la última afirmación: el único solucionador
admisible para trabajadores del tipo 3 es el propio tipo 3, y $\Pi^{2}_{3,3}=0$ daría
$w_3=n_3 z_3/(1+n_3)<z_3$, contradiciendo $w_3\ge z_3$ del Lema~\ref{lem:cotas}.
\end{proof}

\begin{lema}[monotonía del excedente del solucionador]\label{lem:phi}
Sea $s\in[0,1)$ y $\varphi(x)\equiv n(x)\,(s-x)=\dfrac{s-x}{h(1-x)}$ para $x\in[0,s]$. Entonces
$\varphi$ es estrictamente decreciente.
\end{lema}
\begin{proof}
$\varphi(x)=\frac1h\cdot\frac{s-x}{1-x}$ y
$\frac{d}{dx}\frac{s-x}{1-x}=\frac{-(1-x)+(s-x)}{(1-x)^{2}}=\frac{s-1}{(1-x)^{2}}<0$ porque $s<1$.
\end{proof}

\begin{intuicion}
Un solucionador que contrata trabajadores pagándoles su producto individual prefiere trabajadores
\emph{menos} conocedores: al subir el conocimiento del trabajador en $dx$, el excedente por
trabajador cae uno a uno, mientras que el tamaño del equipo solo crece en la proporción $h\,dx$ del
tiempo liberado. Con $h<1$ el primer efecto domina. Este lema es el que decide qué pasa con los más
conocedores.
\end{intuicion}

\begin{lema}[asignación factible]\label{lem:asignacion}
Sea $(w_1,w_2,w_3)$ y $r$ el vector de precios de cualquiera de los Teoremas \ref{th:pre},
\ref{th:aut} o \ref{th:na}. Bajo (A0)--(A3) existe una asignación que satisface $(\Mk)$ y en la que
toda configuración usada tiene beneficio cero.
\end{lema}
\begin{proof}
Asignamos de arriba hacia abajo. En cada paso, ``la configuración que alcanza $w_i$'' significa una
cualquiera de las que realizan el máximo que define $w_i$; existe porque el máximo es sobre un
conjunto finito no vacío, y tiene beneficio cero precisamente por realizarlo.
\begin{enumerate}[label=(\arabic*),leftmargin=2.4em]
  \item \emph{Tipo 3.} Se asigna íntegramente a la configuración que alcanza $w_3$. Esta demanda a lo
  más $n_1\mu_3$ trabajadores del tipo 1, o $n_2\mu_3$ del tipo 2, o $n_a\mu_3$ unidades de compute,
  según cuál sea.
  \item \emph{Tipo 2.} La masa demandada por el paso anterior es a lo más $n_2\mu_3<\mu_2$ por (A2),
  así que está disponible; esos individuos quedan empleados como trabajadores, con beneficio cero por
  la definición de $w_3$. El resto se asigna a la configuración que alcanza $w_2$, que demanda a lo
  más $n_1\mu_2$ trabajadores del tipo 1 si es la de supervisar al tipo 1, y ninguno en los demás
  casos.
  \item \emph{Tipo 1.} La demanda acumulada sobre él es a lo más $n_1(\mu_2+\mu_3)<\mu_1$ por (A1), de
  modo que está disponible, y esas firmas tienen beneficio cero por la definición de $w_2$ y $w_3$.
  Los individuos restantes se asignan a la configuración que alcanza $w_1$: producir por su cuenta si
  $w_1=z_1$, o trabajar bajo un solucionador de IA si el máximo se realiza en esa rama. Esta última no
  está limitada por masas humanas sino por compute.
  \item \emph{Compute.} La demanda total de compute es a lo más $\mu_1/n_1+\mu_2/n_2$ en
  solucionadores de IA más $n_a(\mu_2+\mu_3)$ en trabajadores de IA, cantidad menor que $\mu$ por
  (A3). En el escenario autónomo el compute sobrante se emplea en firmas automatizadas de una capa,
  que con $r=\ai$ ganan cero; no queda compute ocioso y la cláusula final de $(\Mk)$ no se activa. En
  el escenario no autónomo ese compute no puede producir nada, queda ocioso y $(\Mk)$ exige $r=0$, que
  es el valor usado en el Teorema \ref{th:na}.
\end{enumerate}
Todos los humanos quedan empleados y toda configuración usada realiza el máximo que define el salario
del tipo asignado a ella, luego tiene beneficio cero. \qedhere
\end{proof}

\section{Los tres equilibrios}

\begin{teorema}[equilibrio sin IA]\label{th:pre}
Bajo (A1) y (A2) existe un equilibrio sin IA y los salarios son únicos:
\[
  w_1=z_1,\qquad
  w_2=\max\{z_2,\ n_1(z_2-z_1)\},\qquad
  w_3=\max\{z_3,\ n_1(z_3-z_1),\ n_2(z_3-w_2)\}.
\]
\end{teorema}

\begin{proof}
\emph{Unicidad.} Por el Lema~\ref{lem:alcanzado}, $w_1$ es el valor de cero beneficio de una
configuración que emplea al tipo 1: producir solo ($w_1=z_1$) o ser trabajador de un solucionador
$i\in\{2,3\}$ ($w_1=z_i-w_i/n_1$). Si $w_1>z_1$, entonces $\Pi^{1}_1<0$, así que ningún individuo del
tipo 1 produce solo y todos son trabajadores de solucionadores humanos, lo que contradice el
Lema~\ref{lem:masas}. Luego $w_1=z_1$. Para el tipo 2, el Lema~\ref{lem:alcanzado} deja las opciones
$w_2=z_2$, $w_2=n_1(z_2-w_1)=n_1(z_2-z_1)$ y $w_2=z_3-w_3/n_2$; la tercera exigiría que todo el tipo 2
fuese trabajador del tipo 3, excluido por el Lema~\ref{lem:masas}. Por el Lema~\ref{lem:cotas},
$w_2\ge z_2$ y $w_2\ge n_1(z_2-z_1)$, luego $w_2$ es el máximo de ambos. Para el tipo 3, el
Lema~\ref{lem:masas} dice que no es trabajador, así que las opciones son $w_3=z_3$,
$w_3=n_1(z_3-w_1)$ y $w_3=n_2(z_3-w_2)$; el Lema~\ref{lem:cotas} acota $w_3$ por cada una, luego es su
máximo.

\emph{Existencia.} La asignación del Lema \ref{lem:asignacion} satisface $(\Mk)$ con beneficio cero
en las configuraciones usadas. La condición $(\N)$ se cumple: por el Lema \ref{lem:puras} basta
revisar las configuraciones puras, y cada una de ellas aparece como término de alguno de los máximos,
salvo $\Pi^{2}_{i,i}=n_i(z_i-w_i)-w_i$, que es negativo porque $w_i\ge z_i\ge0$.

\emph{Sobre el tipo 2 sobrante.} Si $w_2>z_2$, entonces producir por cuenta propia da pérdidas y todo
el tipo 2 debe emplearse como solucionador de trabajadores del tipo 1, lo que es factible porque
$n_1\mu_2<\mu_1$ por (A1). Si $w_2=z_2$, los sobrantes producen por su cuenta con beneficio cero. En
ambos casos la asignación existe.

\end{proof}

\begin{obs}[quién cobra renta]\label{obs:renta}
El tipo 1 gana exactamente su producto individual: no cobra renta. El tipo 2 cobra renta,
$w_2>z_2$, si y solo si $n_1(z_2-z_1)>z_2$, es decir si y solo si $z_2>\bar z$.
\end{obs}

\begin{teorema}[equilibrio con IA autónoma]\label{th:aut}
Bajo (A0)--(A3), con IA autónoma existe equilibrio, $r^{*}=\ai$ y los salarios son únicos:
\[
  w_1^{*}=\max\Big\{z_1,\ \ai\Big(1-\frac{1}{n_1}\Big)\Big\},\qquad
  w_2^{*}=\max\Big\{z_2,\ \ai\Big(1-\frac{1}{n_2}\Big)\ \text{ó}\ n_a(z_2-\ai),\ n_1(z_2-w_1^{*})\Big\},
\]
\[
  w_3^{*}=\max\big\{z_3,\ n_a(z_3-\ai),\ n_1(z_3-w_1^{*}),\ n_2(z_3-w_2^{*})\big\},
\]
donde en $w_2^{*}$ aparece el término $\ai(1-1/n_2)$ si $z_2\le \ai$ y el término $n_a(z_2-\ai)$ si
$z_2\ge \ai$.
\end{teorema}

\begin{proof}
\emph{Alquiler.} Por (A3) el compute excede todos sus usos con humanos, así que alguna unidad se
emplea en una firma automatizada de una capa; por $(\Z)$, $\Pi^{1}_{AI}=0$, es decir $r^{*}=\ai$.

\emph{Salarios.} El argumento es el del Teorema~\ref{th:pre}, con dos configuraciones más. Para el
tipo 1, el Lema~\ref{lem:alcanzado} da como opciones $z_1$, $\ai-r^{*}/n_1=\ai(1-1/n_1)$ (trabajador
de un solucionador de IA) y $z_i-w_i^{*}/n_1$ (trabajador de un solucionador humano). Si el salario
excediera el máximo de las dos primeras, ningún individuo del tipo 1 podría producir solo ni trabajar
bajo IA sin pérdidas, luego todos serían trabajadores de solucionadores humanos, lo que contradice el
Lema~\ref{lem:masas}. El Lema~\ref{lem:cotas} acota $w_1^{*}$ por las dos primeras, de donde la
fórmula. El tipo 2 se trata igual, usando (A2), y añadiendo la opción de supervisar trabajadores de IA
cuando $z_2\ge\ai$. El tipo 3 no es trabajador (Lema~\ref{lem:masas}), así que su salario es el máximo
de los cuatro valores listados.

\emph{Existencia y $(\N)$.} La asignación del Lema \ref{lem:asignacion} satisface $(\Mk)$, y el
compute sobrante se emplea en firmas automatizadas de una capa, lo que sostiene $r^{*}=\ai$. Para
$(\N)$, el Lema \ref{lem:puras} reduce el problema a las configuraciones puras; cada una es un término
de alguno de los máximos, salvo $\Pi^{2}_{i,i}\le0$, que se sigue de $w_i^{*}\ge z_i\ge0$, y salvo
$\Pi^{tA}_{3}$ y $\Pi^{bA}_{1}$, no admisibles por (A0) excepto en el caso de frontera $z_1=\ai$,
donde $\Pi^{bA}_{1}=-w_1^{*}\le0$.

\end{proof}

\begin{teorema}[equilibrio con IA no autónoma]\label{th:na}
Bajo (A0)--(A3), con IA no autónoma existe equilibrio, $r^{\star}=0$ y los salarios son únicos:
\[
  w_1^{\star}=\max\{z_1,\ \ai\},\qquad
  w_2^{\star}=\max\{z_2,\ \ai\cdot 1\{z_2\le \ai\},\ n_1(z_2-w_1^{\star})\},\qquad
  w_3^{\star}=\max\{z_3,\ n_1(z_3-w_1^{\star}),\ n_2(z_3-w_2^{\star})\}.
\]
\end{teorema}

\begin{proof}
Sin autonomía las firmas de una capa automatizadas y las de trabajadores de IA no producen, así que
todo el compute que se usa va a solucionadores de IA; su demanda es a lo más
$\mu_1/n_1+\mu_2/n_2<\mu$ por (A3), luego sobra compute ocioso y $(\Mk)$ da $r^{\star}=0$. Entonces
$\Pi^{tA}_{j}\le0$ se lee $n_j[\ai-w_j^{\star}]\le0$, es decir $w_j^{\star}\ge \ai$ para los tipos con
$z_j\le \ai$. El resto del argumento de unicidad es idéntico al del Teorema \ref{th:aut}, eliminando
las configuraciones no admisibles, y la existencia es de nuevo el Lema \ref{lem:asignacion}.
\end{proof}

\section{Resultados}

\begin{corolario}[ganadores en el fondo: análogo de la Proposición 5]\label{cor:fondo}
Con IA autónoma,
\[
  w_1^{*}>w_1\iff \ai>\bar z=\frac{z_1}{1-h(1-z_1)} .
\]
Con IA no autónoma, $w_1^{\star}>w_1\iff \ai>z_1=w_1$. Además $\bar z>z_1$ y
\[
  \frac{\bar z}{z_1}=\frac{1}{1-h(1-z_1)}=\frac{n_1}{n_1-1}>1 .
\]
\end{corolario}
\begin{proof}
Por el Teorema~\ref{th:pre}, $w_1=z_1$. Por el Teorema~\ref{th:aut},
$w_1^{*}=\max\{z_1,\ai(1-1/n_1)\}$, luego $w_1^{*}>z_1$ si y solo si $\ai(1-1/n_1)>z_1$; como
$1-1/n_1=1-h(1-z_1)>0$, esto equivale a $\ai>\bar z$. Por el Teorema~\ref{th:na},
$w_1^{\star}=\max\{z_1,\ai\}>z_1$ si y solo si $\ai>z_1$. Finalmente
$1-h(1-z_1)\in(0,1)$ da $\bar z>z_1$, y $1/n_1=h(1-z_1)$ da la expresión del cociente.
\end{proof}

\begin{intuicion}
En ambos regímenes la condición es sobre la \emph{capacidad} de la IA: el trabajador menos conocedor
acepta trabajar bajo un jefe de IA solo si eso le paga más que producir solo. Lo que cambia con la
autonomía es el precio del jefe: con autonomía la IA tiene costo de oportunidad $\ai$ (podría producir
sola) y ese costo se descuenta del salario del trabajador; sin autonomía su costo de oportunidad es
cero y el trabajador se queda con todo el producto del equipo. Por eso el umbral se multiplica por
$n_1/(n_1-1)$: la autonomía no decide \emph{si} hay ganadores abajo, decide qué tan buena debe ser la
IA para que los haya.
\end{intuicion}

\begin{corolario}[el fondo siempre prefiere la IA no autónoma]\label{cor:prefiere}
Si $\ai>z_1$, entonces $w_1^{\star}=\ai>w_1^{*}$.
\end{corolario}
\begin{proof}
$w_1^{\star}=\ai$ y $w_1^{*}=\max\{z_1,\ai(1-1/n_1)\}$. Como $\ai>z_1$ y $\ai(1-1/n_1)<\ai$ porque
$n_1>1$, el máximo es estrictamente menor que $\ai$.
\end{proof}

\begin{corolario}[el tipo automatizado]\label{cor:automatizado}
Supóngase $\ai=z_2$. Entonces $w_2^{*}=w_2^{\star}=z_2$ y, por lo tanto, $w_2^{*}\le w_2$, con
desigualdad estricta si y solo si $z_2>\bar z$. Si además $z_2>\bar z$,
\[
  w_1^{*}-w_1=\frac{w_2-w_2^{*}}{n_1}.
\]
\end{corolario}
\begin{proof}
Con $\ai=z_2$, el término $\ai(1-1/n_2)$ del Teorema~\ref{th:aut} es menor que $z_2$ y
$n_a(z_2-\ai)=0$, así que $w_2^{*}=\max\{z_2,\ n_1(z_2-w_1^{*})\}$. Por el Teorema~\ref{th:aut},
$w_1^{*}\ge \ai(1-1/n_1)=z_2-z_2/n_1$, luego $n_1(z_2-w_1^{*})\le z_2$ y por lo tanto $w_2^{*}=z_2$.
El mismo cálculo con $w_1^{\star}=\ai=z_2$ da $n_1(z_2-w_1^{\star})=0$ y $w_2^{\star}=z_2$. Por la
Observación~\ref{obs:renta}, $w_2=\max\{z_2,n_1(z_2-z_1)\}\ge z_2=w_2^{*}$, con desigualdad estricta
si y solo si $z_2>\bar z$. Para la identidad, si $z_2>\bar z$ entonces $w_2=n_1(z_2-z_1)$,
$w_1^{*}=z_2(1-1/n_1)$ y $w_1=z_1$, de donde
\[
  \frac{w_2-w_2^{*}}{n_1}=\frac{n_1(z_2-z_1)-z_2}{n_1}=z_2-z_1-\frac{z_2}{n_1}
  =z_2\Big(1-\frac{1}{n_1}\Big)-z_1=w_1^{*}-w_1. \qedhere
\]
\end{proof}

\begin{intuicion}
El tipo cuyo conocimiento la IA replica pierde toda su renta, porque compite con un sustituto perfecto
de oferta elástica cuyo precio es $\ai$. La identidad dice adónde va esa renta: se reparte, en partes
iguales, entre los $n_1$ trabajadores del equipo que ese tipo supervisaba. Es el ``efecto share'' del
paper, aquí como identidad contable exacta.
\end{intuicion}

\begin{corolario}[los más conocedores: análogo de la Proposición 6(4)]\label{cor:top}
$w_3^{*}\ge w_3^{\star}$: los más conocedores prefieren, débilmente, la IA autónoma. Si además
$\ai=z_2$, entonces $w_3^{*}\le w_3$: pierden respecto del mundo sin IA.
\end{corolario}
\begin{proof}
Para la primera parte, por los Teoremas~\ref{th:aut} y \ref{th:na} basta comparar término a término.
Como $w_1^{*}\le w_1^{\star}$ y $w_2^{*}\le w_2^{\star}$ (Corolarios~\ref{cor:prefiere} y
\ref{cor:automatizado}), se tiene $n_1(z_3-w_1^{*})\ge n_1(z_3-w_1^{\star})$ y
$n_2(z_3-w_2^{*})\ge n_2(z_3-w_2^{\star})$; el máximo de $w_3^{*}$ contiene además el término
$n_a(z_3-\ai)\ge0$. Luego $w_3^{*}\ge w_3^{\star}$.

Para la segunda, sea $\ai=z_2$, de modo que $w_2^{*}=z_2$ y $n_a(z_3-\ai)=n_2(z_3-z_2)$; los dos
términos centrales del máximo de $w_3^{*}$ coinciden y
$w_3^{*}=\max\{z_3,\ n_2(z_3-z_2),\ n_1(z_3-w_1^{*})\}$. Comparamos con
$w_3=\max\{z_3,\ n_1(z_3-z_1),\ n_2(z_3-w_2)\}$ término a término:
\begin{enumerate}[label=(\roman*),leftmargin=2.2em]
  \item $n_1(z_3-w_1^{*})\le n_1(z_3-z_1)$, porque $w_1^{*}\ge z_1$ por el Teorema~\ref{th:aut};
  \item $n_2(z_3-z_2)<n_1(z_3-z_1)$, por el Lema~\ref{lem:phi} con $s=z_3$ aplicado a $z_1<z_2$.
\end{enumerate}
Ambos términos de $w_3^{*}$ están dominados por términos de $w_3$, y $z_3$ aparece en los dos
máximos, luego $w_3^{*}\le w_3$. \qedhere
\end{proof}

\begin{intuicion}
Los más conocedores pierden por el Lema~\ref{lem:phi}: cambiar trabajadores del tipo 1 por agentes de
IA más capaces encoge el equipo más rápido de lo que sube el excedente por trabajador. Y no lo
compensan por el lado del reparto, porque bajo (A1) sus trabajadores no cobraban renta antes de la IA:
no hay renta ajena que capturar. Aun así prefieren la IA autónoma a la no autónoma, porque sin
autonomía el salario de sus trabajadores es mayor.
\end{intuicion}

\section{La frontera con el modelo del continuo}

El Corolario~\ref{cor:top} dice que los más conocedores pierden, mientras que la Proposición 5 del
paper afirma que siempre ganan. La diferencia está localizada y es demostrable: el resultado del paper
exige $h<h_0$, es decir que \emph{no} existan productores independientes antes de la IA, de modo que
todos los humanos cobren rentas y el efecto \emph{share} opere también arriba. Nuestro supuesto (A1)
produce lo contrario: hay exceso de trabajadores del tipo 1, algunos producen solos y, por el
Teorema~\ref{th:pre}, $w_1=z_1$ exactamente. Sin renta abajo, el jefe no tiene nada que capturar y
solo queda el efecto \emph{match}, negativo por el Lema~\ref{lem:phi}. Es, con otro ropaje, el caso
$h\ge h_0$ que los propios autores identifican como aquel en que los más conocedores pueden perder.

Eliminar (A1) no arregla el problema, lo traslada: sin exceso de ninguna masa el mercado se vacía con
igualdad y el reparto del excedente entre trabajador y solucionador queda indeterminado dentro de un
intervalo. Con un continuo de tipos esa indeterminación desaparece porque la condición de primer orden
$w'(m(z))=n(z)$ selecciona un único reparto en cada punto. Esa es la frontera exacta entre la versión
discreta y la del paper: los tipos finitos bastan para demostrar las condiciones de umbral sobre la
capacidad, pero no para generar rentas en ambos extremos a la vez.

\end{document}
